\subsection{Cross Validation for Picking $\gamma$}
\label{sec:cross}

In this section we show how we used cross validation in our experiments to find $\gamma$.
For each dataset $S$, our framework requires that we 
solve optimization problems of the
form $\min_{\vw} \ell(\vw,S) + \lambda f(\vw, S) + \gamma ||\vw||_2$ for
variable values of $\lambda$,
where $\ell(\vw,S)$ is either MSE (linear regression)
or the logistic regression loss.
For each $\lambda$ we picked $\gamma$ as
a function of this $\lambda$ as follows:

We divided the data into 10 folds (partitions).
Let $S_i$ denote the data of fold $i$ and
$S_{-i}$ denote all the data except fold $i$. 
Let $\Gamma=\{\gamma_1, \ldots, \gamma_k\}$ be a set of $k$ potential values for $\gamma$
we are selecting from and 
$L$ be a vector of size $k$ initialized to all zero.

\begin{algorithm}[h]
 \begin{algorithmic}
% \Procedure{Picking $\gamma$}{} 
\State \For{$j = 1, \ldots, k$ (repeat over values of $\gamma$ in $\Gamma$)}
\State \For{$i = 1, \ldots, 10$ (10-fold cross-validation loop)}
\State Compute $w_i(\gamma_j) \in \arg\min_{w} \ell(w, S_{-i}) + \lambda f(w, S_{-i}) + \gamma_j\|w\|_2$
\State \Comment{train a model $w_i(\gamma_j)$ on all the data except fold $i$}
\State Set $\text{Loss}(i,j) = \ell(w_i(\gamma_j), S_{i}) + \lambda f(w_i(\gamma_j), S_{i})$
\State \Comment{record the test loss of the computed model $w_i(\gamma_j)$ on fold $i$} 
\State $L(j) = L(j) + \text{Loss}(i,j)$.
\State \Comment{add the test loss to the total loss of $\gamma_j$ so far.}
\EndFor
\EndFor
%\EndProcedure
\end{algorithmic}
\caption{Procedure for Picking $\gamma$}
\label{alg:lambda}
\end{algorithm}

 After the loops are over, let $k^*$ denote the entry of $L$ with the smallest value. Then we pick $\gamma_{k^*}\equiv \gamma(\lambda)$ as the regularizer.
So our regularized objective function for this particular $\lambda$ becomes 
\[
\min_w \ell(w,S) + \lambda f(w,S) + \gamma_{\lambda} \|w\|
\]

We observed in our experiments that $\gamma$ increases as $\lambda$ increases. 
Moreover, for any fixed $\lambda$, the value of $\gamma$ chosen for separate models is usually higher than the value
of $\gamma$ chosen for single model.

%
%\subsection*{Step (2): varying $\lambda$ for pareto curves}
%
%For any fixed $\lambda$ we have a regularized objective function derived by the previous step. 
%To derive the pareto curves we do the following. 
%Let $\Lambda  = \{\lambda_1, \ldots, \lambda_k\}$ be the choices for varying values of $\lambda$ and
%$L$ be a vector of size $2 \times k$ initialized to all zero ($L$ saves the error (row 1) and fairness (row 2) losses corresponding
%to each $\lambda$.)
%Again divide the data into 10 partitions. Let $D_i$ denote the partition $i$ of data 
%and $D_{-i}$ denote all the data except partition $i$. 
%
%
%(1) For $j = 1, \ldots, k$ (repeat over values of $\lambda$ in $\Lambda$) \\
%(2)\hspace{5mm}For $i = 1, \ldots, 10$ (10-fold cross-validation loop)\\ 
%(3)\hspace{10mm}Compute the $w_i(\lambda_j) \in \arg\min_{w} \big[\err(w) + \lambda_j F(w) + \gamma_{\lambda_j}\|w\|\big]_{D_{-i}}$ 
%(train a model $w_i(\lambda_j)$ on all the data except partition $i$)\\
%(4)\hspace{10mm}$L(1,j) = L(1,j) + \big[\err(w_i(\lambda_j))\big]_{D_i}$ and $L(2,j) = L(2,j) + \big[F(w_i(\lambda_j))\big]_{D_i}$  (record the test error and test fairness loss of the computed model $w_i(\lambda_j)$  and add it to the total loss so far.)\\
%
%Then we divide $L$ by 10 to get average error and fairness losses for each value of $\lambda$ and then 
%plot the first row of $L$ (error) against it's second row (fairness).